\begin{explanationbox}
	\textbf{\textcolor{CExplanation}{一句话直觉}}
	圆锥曲线的光学性质反映了入射角等于反射角的物理规律：椭圆使从一焦点发出的光线汇聚到另一焦点；双曲线使光线发散，但反向延长线过另一焦点；抛物线使平行于对称轴的光汇聚到焦点，或从焦点发出的光变为平行光。

	\medskip
	\textbf{\textcolor{CExplanation}{核心拆解}}
	\begin{description}[style=nextline, leftmargin=2.8em, labelsep=0.5em, itemsep=0.45em]
		\item[\textbf{要点一（角平分线角色）：}]
			椭圆中，法线平分两焦半径夹角；双曲线中，切线平分两焦半径夹角；抛物线中，切线平分焦半径与平行轴射线的夹角。
		\item[\textbf{要点二（几何基础）：}]
			均源于圆锥曲线的定义及切线性质，可通过解析证明。
	\end{description}

	\medskip
	\textbf{\textcolor{CExplanation}{几何本质（焦点特性）}}
	椭圆定义$|PF_1|+|PF_2|=2a$导致法线平分角；双曲线定义$|PF_1|-|PF_2|=2a$导致切线平分角；抛物线定义$|PF|$等于$P$到准线的距离，导致切线平分角。

	\medskip
	\textbf{\textcolor{CExplanation}{代数意义（解析方法）}}
	通过求导得到切线或法线斜率，利用到角公式可证明角平分关系。代数形式统一为斜率满足的等式。

	\medskip
	\textbf{\textcolor{CExplanation}{考点价值（应用方向）}}
	\begin{itemize}[leftmargin=*, itemsep=0.5em, topsep=0.4em]
		\item \textbf{考法一（角度计算）：} 利用性质直接得到角相等，简化角度计算。
		\item \textbf{考法二（反射路径）：} 在椭圆中，光线经椭圆反射必过另一焦点；在抛物线中，平行光汇聚或光变为平行光，用于求解轨迹。
	\end{itemize}

	\medskip
	\textbf{\textcolor{CExplanation}{顿悟点}}
	“法线夹中点”用于椭圆，“切线夹中点”用于双曲线，“切线分一半”用于抛物线。

	\medskip
	\textbf{\textcolor{CExplanation}{使用场景}}
	\begin{itemize}[leftmargin=*, itemsep=0.5em, topsep=0.4em]
		\item \textbf{场景一（焦点三角形）：} 在椭圆或双曲线的焦点三角形中，由光学性质可建立角度关系。
		\item \textbf{场景二（反射与最值）：} 利用反射性质求解光程最值问题，结合圆锥曲线定义。
	\end{itemize}

\end{explanationbox}
